% Generated draft: requires mathematical and source review.
\documentclass[11pt]{article}
\usepackage[margin=25mm]{geometry}
\usepackage{fontspec}
\setmainfont{Latin Modern Roman}
\usepackage{amsmath,amssymb,mathtools,unicode-math}
\setmathfont{Latin Modern Math}
\usepackage{enumitem,needspace,xurl,hyperref}
\setlength{\parindent}{0pt}
\setlength{\parskip}{5pt}
\setlength{\emergencystretch}{3em}
\newcommand{\N}{\mathbb{N}}
\newcommand{\Z}{\mathbb{Z}}
\newcommand{\Q}{\mathbb{Q}}
\newcommand{\R}{\mathbb{R}}
\newcommand{\C}{\mathbb{C}}
\newcommand{\F}{\mathbb{F}}
\newcommand{\A}{\mathbb{A}}
\newcommand{\PP}{\mathbb{P}}
\newcommand{\E}{\mathbb{E}}
\newcommand{\eps}{\varepsilon}
\newcommand{\dd}{\,\mathrm{d}}
\begin{document}
\section*{UnsolvedMath ID: 30000441}
\textbf{Source problem code:} catalogue entry
\subsection{Definitions and mathematical statement}
\textbf{Analytic descent question.} Does a gauge-equivariant analytic function on flat connections of a flat bundle \(E\to M\) induce an analytic function on the representation variety written in the source as \(\operatorname{Hom}(\pi_1(M),\C)\)? Both gauge equivariance and the displayed target require clarification before this is a determined descent assertion. The group structure on \(\C\) is not specified: its additive group, \(\C^\times\), and a general linear group give different representation spaces. The symbol is preserved rather than silently changed.\par For a rank-\(r\) complex bundle, a conventional setup uses the space \(\mathcal A_{\mathrm{flat}}(E)\) of connections with zero curvature and the gauge group \(\mathcal G(E)\) of bundle automorphisms covering the identity. Choosing a base point and a frame of its fiber yields holonomy in \(\operatorname{Hom}(\pi_1(M),\mathrm{GL}_r(\C))\). Based gauge transformations preserve the framed holonomy; full gauge transformations act by conjugacy. Hence a representation variety, its conjugacy quotient or character variety, and the moduli of flat connections are not interchangeable. A fixed topological bundle may correspond only to the subset of representations with that underlying bundle type. None of these choices is specified in the short source statement.\par Descent as a scalar-valued function requires the original function to be constant on the fibers of the selected holonomy or quotient map: one seeks an analytic \(\widetilde f\) with \(f=\widetilde f\circ\operatorname{Hol}\) on the appropriate domain. Equivariance is not automatically invariance. For example, a rule \(f(g\cdot\nabla)=\chi(g)f(\nabla)\) with a nontrivial scalar character generally describes a section of an associated line object on a quotient rather than an ordinary function, unless the relevant fiber action is trivial or an additional trivialization is specified. For the framed target, invariance along based-gauge fibers is the relevant requirement; for an unframed quotient, the full gauge or conjugacy action must also be treated. The source does not supply the target action, character, line bundle, or trivialization.\par Finally, one must fix whether analytic means real analytic or holomorphic, the topology and analytic category on connections and gauge transformations, and the analytic structure on possibly singular representation spaces. Existence of a set-theoretic quotient function alone does not establish the requested analytic descent. The rank, topology of \(M\), bundle type, and intended function or determinant construction are all missing.\par The supplied triage checked 2026-08-21 labels the problem open and says the report had not constructed the desired analytic object on the representation variety; it also says later determinant work does not establish this exact descent in the extracted generality. Its precise hypotheses are not reproduced. This dated assessment requires independent verification.
\subsection{Short English statement}
Determine when the analytic object on flat connections is constant along the right gauge fibers and descends analytically to the intended representation space. Equivariance may instead give a line-bundle section. The target \(\operatorname{Hom}(\pi_1(M),\C)\), framing, and analytic conventions are unspecified.
\subsection{Sources}
\begin{itemize}\item Local UnsolvedMath registry, permanent ID 30000441, code catalogue entry; \emph{Analytic Functions on Flat-Connection Representation Varieties}, Oberwolfach Reports Open Problems.\item \emph{Analysis and Topology in Interaction (2006)}, Oberwolfach Reports, DOI 10.4171/owr/2006/13, pp. 729--804. \url{https://doi.org/10.4171/owr/2006/13}. Proposed year 2006; no proposer named.\item Supplied literature label \emph{OWR report Analysis and Topology in Interaction}: \url{https://doi.org/10.4171/owr/2006/13}.\item Supplied literature label \emph{renormalized determinants of Dirac operators}: \url{https://msp.org/apde/2022/15-1/apde-v15-n1-p01-s.pdf}. This is a supplied reference only; it was not downloaded.\item Literature triage checked 2026-08-21, status \texttt{open}; record updated 2026-09-16. The \texttt{corrected\_verified} field concerns restored notation and does not define the omitted group or descent action. Verify dated claims independently.\end{itemize}
\end{document}
